\chapter{Implementation Details}
\label{ch:implementation}

This chapter discusses key implementation details, Android-specific considerations, and optimization techniques.

\section{Packet Parsing with Java NIO}

Raw packet parsing uses \code{ByteBuffer} for zero-copy efficiency:

\begin{lstlisting}[language=Kotlin, caption={IPv4 Header Parsing}]
fun parseIpv4(buffer: ByteBuffer): RawPacket? {
    val versionIhl = buffer.get(0).toInt() and 0xFF
    val version = versionIhl shr 4
    if (version != 4) return null
    
    val ihl = (versionIhl and 0x0F) * 4
    val protocol = Protocol.fromNumber(buffer.get(9).toInt() and 0xFF)
    
    val srcIp = formatIpv4(buffer, 12)
    val dstIp = formatIpv4(buffer, 16)
    val (srcPort, dstPort) = extractPorts(buffer, ihl, protocol)
    
    return RawPacket(srcIp, dstIp, srcPort, dstPort, protocol, buffer.remaining())
}
\end{lstlisting}

\section{Token Bucket with Nanosecond Precision}

Token refill uses \code{System.nanoTime()} for precise timing:

\begin{lstlisting}[language=Kotlin, caption={Token Bucket Refill}]
@Synchronized
fun refill() {
    val now = System.nanoTime()
    val elapsed = (now - lastRefillTime) / 1_000_000_000.0
    tokens = minOf(burstBytes.toDouble(), tokens + rateBps * elapsed)
    lastRefillTime = now
}
\end{lstlisting}

\section{Concurrency Management}

\subsection{Thread-Safe Data Structures}

\begin{itemize}
    \item \code{ConcurrentHashMap} for device registry (lock-free reads)
    \item \code{@Synchronized} for token bucket operations
    \item \code{StateFlow} for reactive state propagation
\end{itemize}

\subsection{Coroutine-Based I/O}

Packet processing runs on \code{Dispatchers.IO}:

\begin{lstlisting}[language=Kotlin, caption={Packet Loop Coroutine}]
private val scope = CoroutineScope(Dispatchers.IO + SupervisorJob())

private fun startTunnel() {
    tunInterface = Builder()
        .setSession("QoS Scheduler")
        .addAddress("10.0.0.1", 32)
        .addRoute("0.0.0.0", 0)
        .establish()
    
    scope.launch { runPacketLoop(tunInterface!!) }
}
\end{lstlisting}

\section{Android-Specific Considerations}

\subsection{Foreground Service}

\qos service runs as a foreground service with persistent notification:

\begin{lstlisting}[language=Kotlin, caption={Foreground Service Notification}]
override fun onStartCommand(intent: Intent?, flags: Int, startId: Int): Int {
    createNotificationChannel()
    startForeground(NOTIFICATION_ID, buildNotification())
    return START_STICKY
}
\end{lstlisting}

\subsection{Battery Optimization}

To prevent Android from killing the service:
\begin{itemize}
    \item Foreground service with ongoing notification
    \item \code{START\_STICKY} return value for automatic restart
    \item Request battery optimization exemption (optional)
\end{itemize}

\subsection{VPN Permission}

VPN permission is requested via \code{VpnService.prepare()}:

\begin{lstlisting}[language=Kotlin, caption={VPN Permission Request}]
private fun requestVpnAndStart() {
    val intent = VpnService.prepare(this)
    if (intent != null) {
        vpnPermissionLauncher.launch(intent)
    } else {
        viewModel.startScheduler()
    }
}
\end{lstlisting}

\section{Performance Optimizations}

\subsection{Flow Caching}

Classification results are cached to avoid repeated lookups:

\begin{lstlisting}[language=Kotlin, caption={Flow Cache}]
private val flowCache = ConcurrentHashMap<FlowKey, TrafficCategory>()

fun classify(packet: RawPacket): TrafficCategory {
    val key = FlowKey(packet.srcIp, packet.dstIp, packet.srcPort, packet.dstPort, packet.protocol)
    return flowCache.getOrPut(key) { matchRules(packet.dstPort, packet.protocol) }
}
\end{lstlisting}

\subsection{Conditional Rebalancing}

Rebalancing is triggered only when necessary:

\begin{lstlisting}[language=Kotlin, caption={Conditional Rebalancing}]
private var needsRebalance = false
private var packetCount = 0

fun onPacketProcessed() {
    packetCount++
    if (needsRebalance || packetCount % 1000 == 0) {
        scheduler.rebalanceWithDevices(registry.getAll())
        needsRebalance = false
    }
}
\end{lstlisting}

\section{Error Handling}

All I/O operations are wrapped in \code{runCatching}:

\begin{lstlisting}[language=Kotlin, caption={Safe Packet Forwarding}]
val allowed = scheduler.processPacket(packet)
if (allowed) {
    runCatching { 
        outputStream.write(packet.rawBuffer, 0, packet.length) 
    }.onFailure { 
        Log.e(TAG, "Failed to forward packet", it) 
    }
}
\end{lstlisting}

\section{Testing Strategy}

\subsection{Unit Tests}

\begin{itemize}
    \item Token bucket consume/refill logic
    \item DPI classifier rule matching
    \item Packet parser for IPv4/IPv6
\end{itemize}

\subsection{Integration Tests}

\begin{itemize}
    \item Device registry timeout eviction
    \item Priority persistence across restarts
    \item Rebalancing correctness
\end{itemize}

\subsection{Performance Tests}

\begin{itemize}
    \item Packet processing latency measurement
    \item CPU profiling under 10k pps load
    \item Memory leak detection (24-hour run)
\end{itemize}
